% Numbers sourced from /home/sharaths/projects/game-llm/research/memory/findings.md (F1, F21, F22, F23, F24)
% and /home/sharaths/projects/game-llm/results/exp10_5090/, exp11_seed*, exp12/, exp13_seed*, exp13b/.

Language modeling is reformulated as equilibrium computation across three stages: training dynamics via Magnetic Mirror Descent toward quantal response equilibria, model depth as a weight-tied transformer block solved to a fixed point (EqLM), and multi-model decoding through truthful second-price token auctions. Six hypotheses carrying pre-registered quantitative gates were empirically adjudicated, with every finding subjected to adversarial review. An anytime-unrolled training regime closes the width gap to achieve a weight-tied single-block language model that reaches a ratio of 0.991 of a parameter-matched 12-layer explicit transformer's BLiMP at 121M parameters, trains 2.1 times faster, and gracefully trades off inference cost through a dial over budgets $\{4,8,12\}$ iterations. A trajectory-local contraction penalty yields the first certified 121M-parameter equilibrium (convergence rate 1.0 at 4.0 mean iterations), resolving the gap between quality and certification into two separately solved halves whose naive combination is provably refuted by its own pre-registered prediction. On matrix games, Magnetic Mirror Descent converges linearly (log-linear $R^2$ of 0.995 on matching pennies) where simultaneous gradient play cycles; in preference optimization, the magnetic anchor proves second-order to the preference gradient across four orders of magnitude of regularization strength. Truthful second-price token auctions over domain specialists beat the best single specialist by 23\% and uniform logit averaging by 12\% on mixed-domain perplexity at scoring time, yet this advantage inverts in closed-loop generation, an inversion the adversarial reviewer correctly forced into scope before the follow-up experiment demonstrated its necessity. At equal compute the weight-tied block reaches 0.958 of an explicit transformer's quality with 2.70 times fewer parameters at 46--121M parameters, and a pre-registered twin at one billion parameters on web data locates the boundary of that result: the tied arm fails its kill gate at 1B tokens (perplexity ratio 1.56 against 1.20) and is still closing at 2.5B tokens (1.31), with both arms at chance on public benchmarks. The program closes at that boundary with two formally missed hypotheses, one scale result left open by design, and quality-preserving certification identified as a sharply posed open problem.
